\documentclass[11pt]{article}

\usepackage[margin=1in]{geometry}
\usepackage{times}
\usepackage{setspace}
\usepackage{booktabs}
\usepackage{graphicx}
\usepackage{amsmath}
\usepackage{hyperref}
\usepackage{cite}

\singlespacing
\setlength{\parskip}{4pt}
\setlength{\parindent}{0pt}

% ─── Title ───────────────────────────────────────────────────────────────────
\title{
  \vspace{-1.5em}
  \large\textbf{Re-implementing LoRA: Low-Rank Adaptation of Large Language Models} \\
  \normalsize CS 4/5782 Final Project --- Cornell University, Spring 2026
}
\author{
  Jason Dong (jd876) \quad Dalton Luce (dcl252) \quad Sennet Senadheera (sas639)
}
\date{}

\begin{document}
\maketitle
\vspace{-1em}

% ─── 1. Introduction ─────────────────────────────────────────────────────────
\section*{1. Introduction}

Full fine-tuning of large pre-trained language models requires updating every
parameter, which becomes prohibitively expensive as model sizes grow to billions
of parameters. \textbf{LoRA} (Hu et al., 2021) \cite{lora} addresses this by
freezing the pre-trained weights and injecting trainable low-rank decomposition
matrices into the attention layers. For a weight matrix $W \in \mathbb{R}^{d
\times k}$, LoRA reparameterizes the update as $\Delta W = BA$, where $B \in
\mathbb{R}^{d \times r}$ and $A \in \mathbb{R}^{r \times k}$ with rank $r \ll
\min(d, k)$. This reduces trainable parameters by orders of magnitude with no
additional inference latency, since $W + BA$ can be merged at inference time.

We re-implement LoRA from scratch (without the PEFT library) and evaluate it
against full fine-tuning on RoBERTa-base \cite{roberta}, targeting the results
reported in Table~2 of the original paper.

% ─── 2. Chosen Result ────────────────────────────────────────────────────────
\section*{2. Chosen Result}

We reproduce the RoBERTa-base results from \textbf{Table~2} of Hu et al.\ (2021),
specifically the comparison of full fine-tuning vs.\ LoRA ($r=8$) on two GLUE
\cite{glue} tasks: SST-2 (binary sentiment classification) and MNLI (3-class
natural language inference). The key claim is that LoRA reduces trainable
parameters from $\sim$125M (100\%) to $\sim$0.3M ($\sim$0.3\%) while matching
or exceeding full fine-tuning accuracy (94.8 vs.\ 94.9 on SST-2; 87.6 vs.\ 87.5
on MNLI).

% ─── 3. Methodology ──────────────────────────────────────────────────────────
\section*{3. Methodology}

\textbf{Custom LoRA implementation.}
We implement LoRA entirely from scratch in \texttt{code/lora.py}. The core
class, \texttt{LoRALinear}, wraps any \texttt{nn.Linear} layer with two
additional parameter matrices: $A$ (initialized with kaiming-uniform, equivalent
in scale to the Gaussian initialization specified in the paper) and $B$
(initialized to zero, so $\Delta W = BA = 0$ at the start of training). The
forward pass computes:
\[
  h = W_0 x + \frac{\alpha}{r}(BAx)
\]
where $\alpha$ is a scaling hyperparameter. The function \texttt{apply\_lora()}
freezes all base model parameters, replaces the \texttt{query} and \texttt{value}
projection layers in each of RoBERTa's 12 transformer blocks with
\texttt{LoRALinear}, and keeps the task-specific classification head trainable
(it is randomly initialized per task and must be trained).

\textbf{Hyperparameters.}
We use $r=8$ and $\alpha=16$ (matching Table~2 of the paper), dropout 0.1,
learning rate $2 \times 10^{-4}$ for LoRA and $2 \times 10^{-5}$ for full
fine-tuning, batch size 8, and 2 training epochs.

\textbf{Feasibility modifications.}
To make experiments feasible on a single Colab GPU, we make two adjustments
relative to the paper: (1) we use a random subset of 2,000 training samples
and 500 validation samples per task (vs.\ full datasets of 67K/393K), and (2)
we truncate sequences to a maximum length of 128 tokens rather than RoBERTa's
default of 512, since the vast majority of SST-2 sentences and MNLI sentence
pairs fit within 128 tokens. Both adjustments reduce compute while preserving
the qualitative comparison between methods.

\textbf{Datasets and evaluation.}
SST-2 and MNLI are loaded via HuggingFace \texttt{datasets}. We evaluate
accuracy on \texttt{validation} (SST-2) and \texttt{validation\_matched} (MNLI).

% ─── 4. Results & Analysis ───────────────────────────────────────────────────
\section*{4. Results \& Analysis}

% TODO: Fill in after running experiments in Colab.
% Suggested table format:

\begin{table}[h]
\centering
\small
\begin{tabular}{llrrrr}
  \toprule
  Method & Dataset & Accuracy & Trainable params & \% of total & Train time \\
  \midrule
  Full fine-tuning & SST-2  & 92.4 & 124.6M & 100\%  & 2.0m \\
  LoRA ($r=8$)     & SST-2  & 92.6 & 0.9M & 0.71\% & 1.3m \\
  Full fine-tuning & MNLI   & 79.2 & 124.6M & 100\%  & 1.9m \\
  LoRA ($r=8$)     & MNLI   & 77.2 & 0.9M & 0.71\% & 1.3m \\
  \midrule
  Full fine-tuning & SST-2 (paper) & 94.8 & 125M & 100\% & --- \\
  LoRA ($r=8$)     & SST-2 (paper) & 94.9 & 0.3M & 0.3\% & --- \\
  Full fine-tuning & MNLI  (paper) & 87.6 & 125M & 100\% & --- \\
  LoRA ($r=8$)     & MNLI  (paper) & 87.5 & 0.3M & 0.3\% & --- \\
  \bottomrule
\end{tabular}
\caption{Re-implementation results vs.\ original paper (Table~2). Experiments
run on a Colab GPU with 2,000 training samples, 2 epochs.}
\end{table}

% TODO: Discuss discrepancies, e.g.:
%   - Lower accuracy expected due to smaller training set (2K vs full dataset)
%   - LoRA/baseline gap should still be small, confirming the paper's claim
%   - Parameter counts should exactly match paper values

% ─── 5. Reflections ──────────────────────────────────────────────────────────
\section*{5. Reflections}

% TODO: Fill in after experiments. Suggested talking points:
%   - What worked as expected vs. what surprised you
%   - The importance of keeping the classifier head trainable (easy bug to miss)
%   - Feasibility tradeoffs: smaller subset means lower absolute accuracy
%     but the relative comparison between methods should hold
%   - What you would do differently: train on full dataset, more epochs,
%     ablation over rank r

% ─── References ──────────────────────────────────────────────────────────────
\bibliographystyle{plain}

\begin{thebibliography}{9}

\bibitem{lora}
Hu, E.~J., Shen, Y., Wallis, P., Allen-Zhu, Z., Li, Y., Wang, S., Wang, L.,
\& Chen, W. (2021).
\textit{LoRA: Low-Rank Adaptation of Large Language Models.}
arXiv:2106.09685.

\bibitem{roberta}
Liu, Y., Ott, M., Goyal, N., Du, J., Joshi, M., Chen, D., Levy, O.,
Lewis, M., Zettlemoyer, L., \& Stoyanov, V. (2019).
\textit{RoBERTa: A Robustly Optimized BERT Pretraining Approach.}
arXiv:1907.11692.

\bibitem{glue}
Wang, A., Singh, A., Michael, J., Hill, F., Levy, O., \& Bowman, S. (2018).
\textit{GLUE: A Multi-Task Benchmark and Analysis Platform for Natural Language
Understanding.}
arXiv:1804.07461.

\bibitem{transformers}
Wolf, T., et al.\ (2020).
\textit{HuggingFace's Transformers: State-of-the-art Natural Language Processing.}
arXiv:1910.03771.

\end{thebibliography}

\end{document}
